\begin{lemma}[The polar frame entries of the gradient and Hessian]
Let $n,m\in\mathbb{N}$ with $m\ge1$, let $P\in\mathbb{R}[X_1,\dots,X_n]$ be
homogeneous of degree $m$, let $\mathbb{R}^n$ carry the Euclidean inner product
and let $u(y)=P(y)$. Let $e,v\in\mathbb{R}^n$ and $t\in\mathbb{R}$, and set
\[
  F(s)=u\bigl(\cos(s)e+\sin(s)v\bigr),\qquad
  x=\cos(t)e+\sin(t)v,\qquad
  e_r=x,\qquad e_\theta=-\sin(t)e+\cos(t)v .
\]
Then
\[
  \langle\nabla u(x),e_r\rangle=m\,F(t),
  \qquad
  \langle\nabla u(x),e_\theta\rangle=F'(t),
\]
\[
  D^2u(x)[e_r,e_r]=m(m-1)F(t),
  \qquad
  D^2u(x)[e_r,e_\theta]=(m-1)F'(t),
  \qquad
  D^2u(x)[e_\theta,e_\theta]=F''(t)+m\,F(t).
\]
\end{lemma}
\begin{proof}
Write $\gamma(s)=\cos(s)e+\sin(s)v$, so $F(s)=u(\gamma(s))$, $\gamma(t)=x=e_r$
and $\gamma'(s)=-\sin(s)e+\cos(s)v$, so $\gamma'(t)=e_\theta$. Note $F(t)=u(x)$.

\emph{Step 1: $\langle\nabla u(x),e_r\rangle=m\,F(t)$.} This is the first
conclusion of \noderef{GradientEuler} applied at the point $x$ (with any second
argument), which gives $\langle\nabla u(x),x\rangle=m\,u(x)$; and $e_r=x$,
$u(x)=F(t)$.

\emph{Step 2: $\langle\nabla u(x),e_\theta\rangle=F'(t)$.} By
\noderef{MvPolyFDeriv} applied to $P$ at the point $\gamma(t)=x$, the polynomial
function $z\mapsto P(z)$ on the coordinate space $\mathbb{R}^n$ is Frechet
differentiable at $x$ with derivative $w\mapsto\sum_i(\partial_iP)(x)w_i$.
Composing with the identity map of $\mathbb{R}^n$ read from the Euclidean to the
coordinate norm --- a continuous linear isomorphism, since any two norms on a
finite-dimensional space are equivalent --- the chain rule makes $u$ Frechet
differentiable at $x$ with $Du(x)(w)=\sum_i(\partial_iP)(x)w_i$. By
\noderef{MvPolyGradient} the $i$-th coordinate of $\nabla u(x)$ is
$(\partial_iP)(x)$, so by the coordinate formula for the inner product
$Du(x)(w)=\langle\nabla u(x),w\rangle$ for every $w$. The curve $\gamma$ is
differentiable (each coordinate is a linear combination of $\cos$ and $\sin$, and
a map into a finite-dimensional space is differentiable exactly when its
coordinates are), so by the chain rule $F=u\circ\gamma$ is differentiable at $t$
with
\[
  F'(t)=Du(x)\bigl(\gamma'(t)\bigr)=\langle\nabla u(x),e_\theta\rangle .
\]

\emph{Step 3: $D^2u(x)[e_r,w]=(m-1)\langle\nabla u(x),w\rangle$ for every $w$.}
This is the second conclusion of \noderef{GradientEuler} applied at the point
$x$ with the second argument $w$, together with $e_r=x$.

\emph{Step 4: the two entries involving $e_r$.} Taking $w=e_r$ in Step 3 and
then using Step 1,
\[
  D^2u(x)[e_r,e_r]=(m-1)\langle\nabla u(x),e_r\rangle=(m-1)\,m\,F(t)=m(m-1)F(t).
\]
Taking $w=e_\theta$ in Step 3 and then using Step 2,
\[
  D^2u(x)[e_r,e_\theta]=(m-1)\langle\nabla u(x),e_\theta\rangle=(m-1)F'(t).
\]

\emph{Step 5: the polar--polar entry.} Apply \noderef{GreatCircleSecondDeriv}
with $C=0$, $A=e$, $B=v$ at the parameter $t$. Its curve is
$s\mapsto 0+\cos(s)e+\sin(s)v=\gamma(s)$, so its conclusion reads
\[
  F''(t)=D^2u\bigl(\gamma(t)\bigr)\bigl[\gamma'(t),\gamma'(t)\bigr]
   -\bigl\langle\nabla u(\gamma(t)),\,\cos(t)e+\sin(t)v\bigr\rangle
  =D^2u(x)[e_\theta,e_\theta]-\langle\nabla u(x),e_r\rangle ,
\]
using $\gamma(t)=x$, $\gamma'(t)=e_\theta$ and $\cos(t)e+\sin(t)v=x=e_r$. By
Step 1 the subtracted term is $m\,F(t)$, so rearranging,
\[
  D^2u(x)[e_\theta,e_\theta]=F''(t)+m\,F(t).
\]
Steps 1, 2, 4 and 5 are the five asserted identities.
\end{proof}
